%%=============================================================================
%% Inleiding
%%=============================================================================

\chapter{\IfLanguageName{dutch}{Inleiding}{Introduction}}%
\label{ch:inleiding}

Het correct inschatten van de benodigde personeelsbezetting vormt een belangrijke uitdaging voor kleine detailhandelaars. In sectoren waar de verkoop sterk fluctueert onder invloed van externe factoren zoals weersomstandigheden, weekends, feestdagen en lokale evenementen, is het moeilijk om de personeelsplanning optimaal af te stemmen op de verwachte drukte. Een onderschatting van de personeelsbehoefte kan leiden tot lange wachttijden en overbelasting van medewerkers, wat een negatieve impact heeft op de klanttevredenheid. Een overschatting resulteert dan weer in onnodige loonkosten en een verminderde rendabiliteit.

Deze bachelorproef situeert zich binnen de context van een lokale bakkerij, namelijk Bakkerij Bossu in Nazareth. De focus ligt op het inzetten van operationele data (zoals historische verkoopcijfers en openingstijden) en contextuele data (zoals feestdagen en seizoensinvloeden) in combinatie met machine-learningmodellen om verkoopvoorspellingen te genereren. Op basis van deze voorspellingen wordt vervolgens een berekening gemaakt van de benodigde personeelsbezetting.

Het onderwerp wordt afgebakend tot het ontwikkelen van een proof-of-concept dat verkoopvoorspellingen vertaalt naar een indicatie van het aantal benodigde medewerkers. De effectieve personeelsplanning (zoals individuele uurroosters of contractuele beperkingen) valt buiten de scope van dit onderzoek.

\section{\IfLanguageName{dutch}{Probleemstelling}{Problem Statement}}%
\label{sec:probleemstelling}

Bakkerij Bossu in Nazareth wordt dagelijks geconfronteerd met schommelingen in verkoopvolumes. Deze variaties worden beïnvloed door factoren zoals weekends, feestdagen, seizoensgebonden trends en lokale omstandigheden. Momenteel gebeurt de personeelsplanning hoofdzakelijk op basis van ervaring en intuïtie. Hoewel deze aanpak waardevol is, biedt ze geen objectieve, datagedreven onderbouwing.

Het ontbreken van een voorspellingsmodel kan leiden tot inefficiënties. Bij onderschatting van de drukte ontstaan wachttijden en verhoogde werkdruk voor medewerkers. Bij overschatting worden onnodige loonkosten gemaakt. Dit onderzoek heeft een directe meerwaarde voor Bakkerij Bossu als concrete doelgroep. Daarnaast kan de aanpak relevant zijn voor gelijkaardige lokale bakkerijen of kleine detailhandelszaken met vergelijkbare problematiek.

De primaire doelgroep van deze bachelorproef is de zaakvoerder van Bakkerij Bossu, die verantwoordelijk is voor de personeelsplanning. Secundair kan de ontwikkelde methodologie inspirerend zijn voor andere kleine ondernemingen binnen de retailsector.

\section{\IfLanguageName{dutch}{Onderzoeksvraag}{Research question}}%
\label{sec:onderzoeksvraag}

De centrale onderzoeksvraag van deze bachelorproef luidt als volgt:

\begin{quote}
Hoe kunnen operationele en contextuele data van een lokale bakkerij, in combinatie met machine-learningmodellen, worden ingezet om de benodigde hoeveelheid personeel te bepalen op basis van verkoopvoorspellingen?
\end{quote}

Om deze centrale vraag te beantwoorden, worden volgende deelvragen onderzocht:

\textbf{Deelvragen rond het probleem:}
\begin{itemize}
    \item Hoe sterk variëren de verkoopvolumes tijdens weekends, feestdagen en seizoensgebonden periodes?
    \item Welke patronen en trends zijn zichtbaar in de historische verkoopdata?
    \item Welke operationele en contextuele variabelen hebben een significante invloed op de verkoop?
\end{itemize}

\textbf{Deelvragen rond de oplossing:}
\begin{itemize}
    \item Welke machine-learningmodellen zijn het meest geschikt voor het voorspellen van verkoopcijfers in een retailcontext?
    \item Hoe presteren verschillende regressiemodellen en ensemble-methoden in termen van nauwkeurigheid?
    \item Hoe kunnen verkoopvoorspellingen vertaald worden naar een concrete berekening van de benodigde personeelsbezetting op basis van vooraf gedefinieerde capaciteitsregels?
\end{itemize}

\section{\IfLanguageName{dutch}{Onderzoeksdoelstelling}{Research objective}}%
\label{sec:onderzoeksdoelstelling}

Het doel van deze bachelorproef is het ontwikkelen van een proof-of-concept dat Bakkerij Bossu ondersteunt bij het plannen van personeel op basis van datagedreven verkoopvoorspellingen.

Concreet omvat dit:

\begin{itemize}
    \item Het analyseren van historische operationele en contextuele data.
    \item Het trainen, vergelijken en evalueren van meerdere machine-learningmodellen voor verkoopvoorspelling.
    \item Het selecteren van het best presterende model op basis van objectieve evaluatiemetrics zoals RMSE en MAE.
    \item Het ontwikkelen van een prototype dat voorspelde verkoopcijfers omzet in een berekening van het benodigde aantal medewerkers, op basis van vooraf gedefinieerde capaciteitsregels.
    \item Het formuleren van aanbevelingen voor de praktische implementatie binnen de bakkerij.
\end{itemize}

De bachelorproef is succesvol wanneer een werkend proof-of-concept wordt opgeleverd dat betrouwbare verkoopvoorspellingen genereert en een onderbouwde indicatie geeft van de benodigde personeelsbezetting.

\section{\IfLanguageName{dutch}{Opzet van deze bachelorproef}{Structure of this bachelor thesis}}%
\label{sec:opzet-bachelorproef}

De rest van deze bachelorproef is als volgt opgebouwd:

In Hoofdstuk~\ref{ch:stand-van-zaken} wordt een overzicht gegeven van de stand van zaken binnen het onderzoeksdomein, op basis van een literatuurstudie rond verkoopvoorspelling, machine learning en personeelsplanning.

In Hoofdstuk~\ref{ch:methodologie} wordt de methodologie toegelicht en worden de gebruikte onderzoekstechnieken besproken om een antwoord te kunnen formuleren op de onderzoeksvragen.

In Hoofdstuk~\ref{ch:resultaten} worden de resultaten van de data-analyse en de modelvergelijking besproken.

In Hoofdstuk~\ref{ch:proof-of-concept} wordt het ontwikkelde proof-of-concept toegelicht en wordt beschreven hoe de verkoopvoorspellingen vertaald worden naar personeelsbehoefte.

In Hoofdstuk~\ref{ch:conclusie}, tenslotte, wordt de conclusie gegeven en een antwoord geformuleerd op de onderzoeksvragen. Daarbij wordt ook een aanzet gegeven voor toekomstig onderzoek binnen dit domein.